\chapter{CONCLUSIONS ET PERSPECTIVES}
\thispagestyle{empty}

\section{Conclusions}

Ce projet a developpe et analyse un systeme de recommandation d'articles
Wikipedia fonde sur les graphes et centre sur l'espace Custom-WikiCS. A partir
de la topologie Wiki-CS, le travail a enrichi les noeuds avec le texte reel de
Wikipedia, genere des embeddings semantiques et structurels, compare plusieurs
familles de recommandation et expose les resultats a travers une application
prototype locale.

Plusieurs conclusions peuvent maintenant etre formulees clairement.

Premiere conclusion, le graphe nettoye d'articles contient un signal structurel
et semantique significatif. Le comportement en degre a queue lourde,
l'assortativite semantique et une structure communautaire non triviale
justifient tous l'usage de methodes de recommandation sensibles au graphe.

Deuxieme conclusion, le texte et la structure sont tous deux utiles, mais ne se
comportent pas de maniere identique. La recuperation semantique seule reste
forte et interpretable. La recuperation structurelle seule demeure puissante,
en particulier sous des metriques larges de discrimination des aretes.
Cependant, la famille de recommandation la plus forte dans le protocole
d'evaluation corrige est la famille hybride corrigee.

Troisieme conclusion, la correction methodologique ulterieure constitue l'une
des contributions les plus importantes du travail. Le protocole revise
d'evaluation structurelle non orientee a modifie l'interpretation de la
comparaison initiale a huit modeles et a montre que le recit anterieur de
domination de Node2Vec n'etait pas assez robuste pour rester la conclusion
finale.

Quatrieme conclusion, le projet a clarifie la difference entre force de
classement et qualite de recommandation. La famille metrique additionnelle ont montre qu'un modele peut etre bon pour
la discrimination de liens tout en se comportant differemment comme
recommandeur en liste courte.

Enfin, l'application locale a montre que les resultats experimentaux peuvent
etre transformes en une interface exploitable pour la recherche, la selection
de modele, l'inspection des recommandations et une presentation optionnelle en
parcours d'apprentissage. Elle a egalement montre qu'un organisateur LLM
contraint, complete par un repli heuristique, constitue une solution plus
robuste qu'une generation libre pour cette couche de presentation.

\section{Direction finale recommandee pour le modele}

Si l'objectif est la cible applicative reelle d'une recommandation d'articles
en tete de classement, la famille la plus defendable comme choix principal est
la famille hybride corrigee. Au sein de cette famille:
\begin{itemize}
    \item \textit{BGE pooled + Node2Vec (corrected)} constitue un choix fort
    lorsque Hit@10 est prioritaire;
    \item \textit{BGE-M3 + Node2Vec (corrected)} constitue un choix fort
    lorsque MRR, Recall, MAP et plusieurs metriques secondaires de classement
    sont prioritaires.
\end{itemize}

Cette recommandation reflete le cadre final d'evaluation corrigé.

\section{Limites}

Le projet presente egalement des limites claires:
\begin{itemize}
    \item les experiences restent restreintes au sous-domaine informatique de
    Wikipedia;
    \item la base de code demeure fortement centree sur les notebooks et n'a
    pas ete completement refactorisee en modules reutilisables;
    \item l'organisateur par LLM reste une couche de presentation plutot qu'un
    moteur de recommandation valide;
    \item les contraintes de materiel local limitent l'echelle de
    l'entrainement et l'etendue de l'exploration de modeles.
\end{itemize}

\section{Travaux futurs}

Les directions futures les plus utiles sont:
\begin{itemize}
    \item etendre le cadre de recommandation au-dela du sous-domaine actuel tout
    en preservant la proprete de l'evaluation;
    \item tester des strategies de fusion plus avancees que la concatenation
    brute ou le pooling fixe;
    \item ameliorer la couche applicative avec une analyse qualitative plus
    riche et une presentation plus stable des parcours d'apprentissage;
    \item etudier si les metriques de qualite de recommandation peuvent etre
    integrees directement dans la selection de modele plutot que seulement dans
    une analyse a posteriori.
\end{itemize}

Dans l'ensemble, le projet aboutit a un resultat empirique coherent: pour cette
tache de recommandation d'articles Wikipedia, la meilleure direction finale
n'est ni un modele purement semantique ni un modele purement structurel, mais
un systeme hybride soigneusement controle, evalue sous un protocole oriente
recommandation et exempt de fuite.
